\documentclass{rrxiv}
\rrxivid{rrxiv:2605.00003}
\rrxivversion{v2}
\rrxivprotocolversion{0.1.0}
\rrxivlicense{CC-BY-4.0}
\rrxivtopics{stat.ML,cs.LG}

\title{Reproducibility budgets for ML preprints}
\author{Example Author A \and Example Author B \and Example Author C}
\date{2026-05-12}

\begin{document}
\maketitle

\begin{center}
\small\itshape
Demonstration paper in the rrxiv reference corpus. The canonical machine-readable version lives at \href{https://rrxiv.com/papers/rrxiv:2605.00003}{rrxiv.com/papers/rrxiv:2605.00003}.
\end{center}

\begin{abstract}
We argue that preprint platforms should require authors to declare a reproducibility budget at submission time — an explicit specification of the compute, data, and human-time required to replicate the paper's central claims. We evaluate adoption costs over 12 months and report mixed results: smaller labs benefit, large labs resist.
\end{abstract}

\section{Introduction}
We argue that preprint platforms should require authors to declare a reproducibility budget at submission time — an explicit specification of the compute, data, and human-time required to replicate the paper's central claims. We evaluate adoption costs over 12 months and report mixed results: smaller labs benefit, large labs resist.

This document is a structured encoding of the paper in the \texttt{rrxiv} protocol's Canonical Intermediate Representation (CIR). It engages with the topics \texttt{stat.ML} and \texttt{cs.LG}. The encoding registers 3 formal claims (1 contradicted, 1 partial, 1 replicated). Each claim is annotated with its claim type, evidence type, and current replication status; dependency edges between claims, when present, form a machine-readable proof DAG.

\section{Methodology}
We follow the \texttt{rrxiv} convention of separating \emph{claims} (the proposition under consideration) from \emph{evidence} (the argument or data supporting it). Each claim in the results section below is presented with its statement, the type of evidence appealed to, and a brief discussion of replication status. Where claims depend on prior results --- internal or external --- the dependency is recorded in the CIR as a \texttt{\textbackslash dependson} edge, so the full inferential structure is machine-traversable. Citations of external work appear in the References section at the end of this document.

\section{Results: registered claims}
\subsection*{Claim 1}
\begin{claim}[Claim 1]
\label{claim:c1}
Reproducibility-budget declarations correlate with 2.3x higher replication success rates (p \textless{} 0.05).

\emph{Replication status: replicated.}
\end{claim}
This claim is an empirical observation supported by data. As of the encoding date, it has been independently replicated.

\subsection*{Claim 2}
\begin{claim}[Claim 2]
\label{claim:c2}
Mandatory budget declarations reduce submission volume by 15\% in the first quarter.

\emph{Replication status: contradicted.}
\end{claim}
This claim is an empirical observation supported by data. As of the encoding date, it has been contradicted by subsequent work.

\subsection*{Claim 3}
\begin{claim}[Claim 3]
\label{claim:c3}
Budgets self-correct via community feedback over a 6-month window.

\emph{Replication status: partial.}
\end{claim}
This claim is an empirical observation supported by data.

\section{Discussion}
The claim graph above is the primary product of this paper. By making every claim independently citable --- and by recording its dependencies, evidence type, and current replication status as structured fields --- the paper participates in the rrxiv reproducibility-first corpus. Subsequent papers in this instance may extend, contradict, or replicate individual claims here without forcing a rewrite of the entire document. See the canonical version online for the live discourse layer.

\section{References}
\begin{itemize}[leftmargin=*]
\item Computational reproducibility at scale
\item Reproducibility in machine learning
\end{itemize}
\end{document}
